% =====================================================================
\section{Справочник}
% =====================================================================
\subsection*{«Опасные» выводы ESP32-S3}

\begin{center}\small
\begin{tabularx}{\textwidth}{@{}l X@{}}
\toprule
Выводы & Особенность \\
\midrule
GPIO0, 3, 45, 46 & strapping: влияют на режим загрузки; GPIO0 --- кнопка BOOT \\
GPIO19, 20 & нативный USB (D$-$, D$+$) \\
GPIO26--32 & SPI-Flash --- не использовать \\
GPIO33--37 & заняты при Octal Flash/PSRAM (35--37 --- при Octal PSRAM) \\
GPIO39--42 & JTAG (MTCK, MTDO, MTDI, MTMS) --- свободны, если не нужна JTAG-отладка по этим выводам \\
GPIO43, 44 & UART0 (TX, RX) --- сообщения загрузчика \\
GPIO11--20 & ADC2 --- недоступен при работающем Wi-Fi \\
GPIO22--25 & не существуют \\
\bottomrule
\end{tabularx}
\end{center}

\subsection*{Что изучать дальше}

\begin{figure}[H]
\centering
\begin{tikzpicture}[node distance=5mm and 6mm, b/.style={blk,minimum width=28mm}]
  \node[b,fill=blkcpu] (base) {Базовый курс\\ (уроки 0--9)};
  \node[b,fill=blkrf,above right=4mm and 12mm of base] (ble) {Bluetooth LE\\ \scriptsize NimBLE};
  \node[b,fill=blkrf,right=12mm of base] (mqtt) {MQTT, HTTPS,\\ OTA-обновления};
  \node[b,fill=blkrf,below right=4mm and 12mm of base] (usb) {USB HID/MIDI\\ \scriptsize TinyUSB};
  \node[b,fill=blkper,right=of ble] (cam) {Камера,\\ LCD-интерфейс};
  \node[b,fill=blkper,right=of mqtt] (ai) {ESP-DL,\\ TFLite Micro};
  \node[b,fill=blkper,right=of usb] (idf) {Чистый ESP-IDF\\ \scriptsize CMake, menuconfig};
  \foreach \t in {ble,mqtt,usb} \draw[arr] (base.east) -- (\t.west);
  \draw[arr] (ble) -- (cam); \draw[arr] (mqtt) -- (ai); \draw[arr] (usb) -- (idf);
\end{tikzpicture}
\caption{Возможные направления для дальнейшего изучения. Для лампы естественное продолжение ---
MQTT (интеграция с умным домом) и OTA (обновление прошивки по воздуху)}
\end{figure}

\subsection*{Полезные источники}
\begin{itemize}
  \item ESP32-S3 Datasheet и Technical Reference Manual --- \url{https://www.espressif.com/en/support/documents/technical-documents}
  \item Документация Arduino-ESP32 --- \url{https://docs.espressif.com/projects/arduino-esp32/}
  \item Руководство ESP-IDF --- \url{https://docs.espressif.com/projects/esp-idf/en/latest/esp32s3/}
  \item Описание платы DevKitC-1 --- раздел \emph{ESP32-S3-DevKitC-1 User Guide} на \url{https://docs.espressif.com}
\end{itemize}
